% theory_iec60909.tex —— 短路计算 通用理论：IEC 60909 等效电压源法
% 本文件为理论层章节，遵循 docs/modules/THEORY_WRITING_GUIDE.md。

\section{IEC 60909 等效电压源法理论}\label{sec:sc-theory-iec}

\begin{theorynote}
本章为通用数学理论，按 IEC 60909-0 标准体系与经典教材撰写，覆盖等效电压源法的
方法论与叠加原理视角、短路电流时域分解与峰值系数 $\kappa$ 的物理来源、同步电机
电抗谱系（$x_d''/x_d'/x_d$）的电机学解释、修正系数体系（$K_G/K_T/K_S$）、开断电流
衰减因子 $\mu$、稳态电流系数 $\lambda_{\max}/\lambda_{\min}$、热等效系数 $m/n$，
以及近端/远端短路分类。与平台实现（\srcpath{src/short\_circuit/}）的对应关系见
本章末表；表内只区分本模块实现与其他专用物理模块归属，不保留待实现占位。
\end{theorynote}

\subsection{记号与约定}
\begin{itemize}
  \item $U_n$：系统标称线电压；$c$：电压因子（$c_{\max}$ 用于最大、$c_{\min}$ 用于
        最小短路电流计算）；
  \item $I_k''$：初始对称短路电流（有效值）；$i_p$：峰值；$I_b$：对称开断电流；
        $I_k$：稳态短路电流；$I_{th}$：热等效短路电流；
  \item $Z_k$：故障点等效短路阻抗；$R/X$：其电阻-电抗比；$\tau=X/(\omega R)$：
        直流分量衰减时间常数；
  \item $x_d'',x_d',x_d$：同步电机直轴次暂态/暂态/稳态电抗（电机额定基准标幺）；
        $I_{rG}$：发电机额定电流；$t_{\min}$：最短开断延时；$T_k$：短路持续历时。
\end{itemize}

\subsection{等效电压源法与叠加原理}
\begin{definition}[等效电压源法]
IEC 60909-0~\cite{iec:iec60909-0} 将故障前网络在故障点处置换为\emph{唯一}有源元件——等效电压源
$cU_n/\sqrt3$（相电压），其余所有电源均以内阻抗表示、非旋转负荷与非故障相关支路
（线路电容、并联导纳）忽略，于是
\begin{equation}
I_k''=\frac{cU_n}{\sqrt3\,Z_k},
\end{equation}
$Z_k$ 为从故障点看入的等效短路阻抗（不平衡故障按第~\ref{sec:sc-theory-symcomp}~章
的序网组合）。
\end{definition}

\begin{remark}[为什么用 $cU_n/\sqrt3$ 而不是故障前实测电压]
等效电压源是\emph{计算约定}而非物理电源。因子 $c$ 一次性计入：
(i) 系统电压在时间与空间上的分布（运行电压可高于 $U_n$）；
(ii) 变压器分接头位置的离散变化；
(iii) 忽略负荷电流与潮流工况带来的误差；
(iv) 电机次暂态行为的等效折让。
因此同一网络的最大/最小短路电流计算分别采用 $c_{\max}/c_{\min}$，对应设备校核与
保护灵敏性两类口径。
\end{remark}

\begin{remark}[叠加原理视角]
故障可等效为在故障点注入与故障前电压 $V_{pre}$ 等值反向的补偿电压源，与故障前
状态叠加。网络线性时，故障分量只由该补偿源驱动：把全部内电势置零、仅保留故障点
源，即得等效电压源法的电路结构。故障前负荷电流分量被忽略，其影响由 $c$ 折中
补偿；这正是该方法相对潮流初值驱动精确法的近似所在。
\end{remark}

\begin{definition}[近端与远端短路]
\emph{近端短路}：至少一台同步电机贡献的初始短路电流超过其两倍额定电流
（$I_{kG}''>2I_{rG}$），或双侧电源各自贡献均占主导——交流分量随时间显著衰减，
需计及 $\mu$、$\lambda$ 等修正。\emph{远端短路}：交流分量幅值视为不变，
$I_k=I_b=I_k''$。分类决定后续特征电流的计算路径。
\end{definition}

\subsection{短路电流的时域分解与峰值系数}
考察 $R$--$L$ 串联支路在 $t=0$ 合闸于正弦电压 $u(t)=\sqrt2\,U\sin(\omega t+\psi)$：
\begin{equation}
i(t)=\sqrt2\,I_k''\Bigl[\sin(\omega t+\psi-\theta)-\sin(\psi-\theta)\,e^{-t/\tau}\Bigr],
\qquad \theta=\arctan\frac{\omega L}{R},\quad \tau=\frac{L}{R}=\frac{X}{\omega R}.
\end{equation}
第一项为强制交流分量，第二项为衰减的非周期（直流）分量，最坏情形
$|\sin(\psi-\theta)|=1$。首个峰值约在半周波 $t\approx0.01$~s（50~Hz）处：
\begin{equation}
i_p=\kappa\sqrt2\,I_k'',\qquad
\kappa=1+e^{-0.01/\tau}=1+e^{-\pi R/X}\ \ (f=50~\mathrm{Hz}).
\end{equation}
IEC 60909-0 采用工程拟合式
\begin{equation}
\kappa=1.02+0.98\,e^{-3R/X},
\end{equation}
在常用 $R/X$ 范围内与上式偏差约 $1\%$；极限行为一致：$R/X\to0$ 时 $\kappa\to2$
（全偏移），$R/X\to\infty$ 时 $\kappa\to1.02\approx1$（无直流分量）。

\begin{remark}[网状网络的 $\kappa$ 方法 A/B/C]
网状网络中各馈入支路 $R/X$ 不同，标准给出三种等效方法：方法 A（对故障电压等级中
承载分量短路电流的全部支路及馈入该电压等级的变压器取最小 $R/X$，即最大时间常数）、
方法 B（取故障点等效阻抗的 $R/X$，并对结果乘 $1.15$ 安全修正，网状时
$\kappa$ 上限 $1.8$）、方法 C（等效频率法）。方法 C 在 50/60 Hz 系统分别取
$f_c=20/24$ Hz，即 $f_c/f=0.4$；将无源支路与电压源内阻的电抗乘 $0.4$，同步发电机
电阻改用 $R_{Gf}/X_d''=0.15$（低压）、$0.07$（高压且 $S_r<100$ MVA）或 $0.05$
（高压且 $S_r\ge100$ MVA），求频率修正网络的驱动点阻抗 $Z_{c,kk}$，再取
\begin{equation}
  \left(\frac RX\right)_C=\frac{\Re Z_{c,kk}}{\Im Z_{c,kk}}\frac{f_c}{f},\qquad
  \kappa_C=1.02+0.98\exp\!\left[-3(R/X)_C\right].
\end{equation}
式 (59) 最终写成 $i_p=\sqrt2(\kappa I_{k1}''+I_{k2}'')$；静止发电机和跟网换流器
作为限流电流源计入 $I_{k2}''$，不获得非周期分量放大。
\end{remark}

\begin{theorynote}
实现完整区分 A/B/C：方法 A 由故障电压解识别实际承载分量电流的支路，并把电压源内阻和
馈入变压器一并纳入最小 $R/X$，对全部电压源贡献使用同一 $\kappa_A$；方法 B 且网状时取
$\min(1.8,1.15\kappa)$；方法 C 建立独立频率修正稀疏因子。OpenDSS oracle 也以外部支路电流
独立识别同一集合，不能用实现返回的 $\kappa$ 自证。
\end{theorynote}

\subsection{同步电机电抗谱系的电机学解释}
故障瞬间绕组磁链守恒，转子阻尼绕组与励磁绕组依次对定子电枢反应形成屏蔽，使
电机对外呈现的电抗随时间逐档放大：
\begin{equation}
x_d''<x_d'<x_d,
\end{equation}
$x_d''$（阻尼+励磁屏蔽，衰减时间常数 $T_d''$，数十 ms）、$x_d'$（励磁屏蔽，
$T_d'$，数百 ms 至秒级）、$x_d$（稳态）。近端故障机端电流包络为
\begin{equation}
i_{ac}(t)=\sqrt2\Bigl[(I''-I')e^{-t/T_d''}+(I'-I)e^{-t/T_d'}+I\Bigr],
\qquad I''=\frac{E_q}{x_d''},\ I'=\frac{E_q'}{x_d'},\ I=\frac{E_f}{x_d}.
\end{equation}
IEC 60909 的方法论是\emph{避免逐机时域仿真}：用 $x_d''$ 计算 $I_k''$，再以系数
$\mu$（开断时刻）、$\lambda$（稳态）、$\kappa$（峰值）折算后续时段，把上式的
衰减谱折叠进少量代数系数。

\begin{theorynote}
IEC 准静态路径以 $x_d''$ 计算初始量，并按标准的 $\mu$、$q$、$\lambda$ 代数规则计算开断与
稳态量；它不以 $x_d'$、$T_d'$、$T_d''$ 做时域积分。这是 IEC 60909 计算法本身的模型选择，
不是本实现用 $x_d$ 序网猜测稳态量；动态包络由暂态模块承担。
\end{theorynote}

\subsection{修正系数体系}
IEC 60909 对电源设备阻抗引入修正系数，使等效电压源口径下各贡献的幅值/相位
更接近真机行为：
\begin{itemize}
  \item \textbf{发电机 $K_G$}（直接并网机组）：
        \begin{equation}
        Z_{G,K}=K_G Z_G,\qquad
        K_G=\frac{U_n}{U_{rG}(1+p_G)}
            \frac{c_{\max}}{1+x_d''\sin\varphi_{rG}},
        \end{equation}
        折算母线/铭牌电压基准、永久电压偏差 $p_G$ 与额定功率因数下内电势差异；
  \item \textbf{变压器 $K_T$}（网络两/三绕组变压器）：
        \begin{equation}
        Z_{T,K}=K_T Z_T,\qquad K_T=\frac{0.95\,c}{1+0.6\,x_T},
        \end{equation}
        $x_T$ 取铭牌基准标幺值；计入分接头离散与饱和引起的阻抗不确定度；
  \item \textbf{电站机组 $K_S$}：发电机与升压变作为整体从高压侧等效，
        修正含额定变比平方项与 $|x_d''-x_T|$ 的组合项，适用于机组建模到
        升压变高压侧的场合。
\end{itemize}
实现组装阻抗修正始终使用设备所在电压等级的 $c_{\max}$；本次计算的等效电压源才按
最大/最小工况使用 $c_{\max}/c_{\min}$。这一区分避免最小短路把设备修正系数错误缩小。

\subsection{开断电流衰减因子 \texorpdfstring{$\mu$}{μ} 与电动机系数 \texorpdfstring{$q$}{q}}
\begin{definition}[对称开断电流与 $\mu$、$q$]
开断电流 $I_b$ 为开关触头分离时刻（最短延时 $t_{\min}$ 后）的对称短路电流有效值。
近端同步机贡献按 $I_{bG}=\mu I_{kG}''$ 衰减，$\mu$ 是 $t_{\min}$ 与电流比
$I_{kG}''/I_{rG}$ 的二元函数（标准以曲线与分段解析式给出，$\mu\le1$）；
异步电动机贡献按 $I_{bM}=\mu q I_{kM}''$，$q$ 随每对极额定有功功率与 $t_{\min}$
增大而减小（小电机衰减更快）。远端短路 $\mu=1$，$I_b=I_k''$。
\end{definition}

\subsection{稳态电流系数 \texorpdfstring{$\lambda_{\max}/\lambda_{\min}$}{λmax/λmin}}
\begin{definition}[$\lambda$ 系数]
近端短路的稳态电流由励磁响应主导，标准以额定电流倍数给出：
\begin{equation}
I_{kG}=\lambda\,I_{rG},
\end{equation}
$\lambda_{\max}$ 用于最大稳态（校核过励/持续热稳定），$\lambda_{\min}$ 用于最小
稳态（校核保护灵敏性）；二者来自制造商数据或按标准图 15/16 查取，不能由 $x_d/x_q$
臆造。远端短路的网络馈入保持 $I_k=I_k''$；异步电机机端三相故障的稳态贡献为零；多馈近端
故障按第 11.2.7 条取 $I_k=I_{bmo}$，其中 $I_{bmo}$ 是删除全部异步电动机贡献后按公式 (77)
计算的开断电流。端子馈电静态励磁机组发生机端故障时，最大与最小计算均取
$\lambda_{\max}=\lambda_{\min}$，而不是取零。
\end{definition}

\subsection{热等效系数 \texorpdfstring{$m$}{m} 与 \texorpdfstring{$n$}{n}}
\begin{definition}[热等效短路电流]
\begin{equation}
I_{th}=I_k''\sqrt{m+n},
\end{equation}
$m$、$n$ 分别为非周期分量与交流分量衰减的热效应系数。由直流分量
$i_{dc}(t)=\sqrt2 I_k''(\kappa-1)e^{-t/\tau}$ 的能量积分定义：
\begin{equation}
m=\frac{1}{T_k\,(I_k'')^2}\int_0^{T_k} i_{dc}^2\,dt
=(\kappa-1)^2\,\frac{\tau}{T_k}\bigl(1-e^{-2T_k/\tau}\bigr),
\end{equation}
标准附录 A 给出可编程解析式
\begin{equation}
m=\frac{\exp\!\left(4fT_k\ln(\kappa-1)\right)-1}
        {2fT_k\ln(\kappa-1)},
\end{equation}
其连续极限为 $m(1)=0$、$m(2)=2$。令 $q=I_k''/I_k$、
$q'=q/(0.88+0.17q)$、$T_d'=3.1\,\mathrm{s}/q'$，则 $q=1$ 时 $n=1$，
$q\ge1.25$ 时按附录 A 六项闭式计算；$1<q<1.25$ 按图 19 在端点之间插值。
\end{definition}

\begin{theorynote}
\srcpath{calculate\_thermal\_factors\_sc} 逐式实现附录 A，并以 \texttt{expm1} 处理
$\kappa\to2$ 极限、以解析极限处理 $I_k\to0$；结果中的 \fld{thermal\_m/thermal\_n}
允许调用方审计 $I_{th}$ 的组成。
\end{theorynote}

\subsection{与实现的对应}
\begin{center}\footnotesize
\begin{longtable}{@{}L{6.8cm}L{8.6cm}@{}}
\toprule
理论条目 & 实现状态\\
\midrule
\endfirsthead
\toprule
理论条目 & 实现状态\\
\midrule
\endhead
等效电压源与 $c_{\max}/c_{\min}$ 电压因子 &
\implfull{include/hacdcpf/analysis/short\_circuit.hpp:get\_voltage\_factor\_sc}（实现取值约化见第~\ref{sec:sc-overview}~章）\\
故障点三序戴维南等效（选择性求逆 $Y_sz_k=e_k$） &
\implfull{src/short\_circuit/short\_circuit.cpp:build\_detailed\_sparse\_context}\\
$\kappa$ 工程式与等效电压源/电流源分解（IEC 60909-0 式 59） &
\textbf{已实现}\quad \srcpath{src/short\_circuit/short\_circuit.cpp:calculate\_kappa\_basic\_sc}；A/B/C 对电压源等值使用单一方法系数，静止发电机和跟网换流器取 $\kappa=1$\\
网状网络 $\kappa$ 方法 A/B/C 完整体系 &
\textbf{已实现}\quad A 为参与支路统一最小 $R/X$；B 为故障点等值与网状 1.15 修正；C 为独立 $f_c/f=0.4$ 稀疏等值与发电机 $R_{Gf}$\\
近端/远端显式分类 &
\textbf{已实现}\quad 以单机贡献 $I_{kG}^{\prime\prime}>2I_{rG}$ 判定近端，并区分单馈、多馈、网络馈入与异步电机稳态规则\\
电抗谱系 $x_d''/x_d'/x_d$ 与时间常数 &
\textbf{已实现}\quad IEC 准静态路径按 $x_d^{\prime\prime}+\mu/q/\lambda$ 代数体系执行；时域 $x_d'/T_d'/T_d^{\prime\prime}$ 包络不属于本入口\\
修正系数 $K_G/K_T/K_S$ &
\textbf{已实现}\quad \srcpath{generator\_kg}、\srcpath{transformer\_kt} 与 \srcpath{build\_iec\_network\_corrections}；$K_S$ 区分 OLTC/非 OLTC，并对机组内部故障切换到内部 $K_G$\\
开断衰减 $\mu$、电动机 $q$ 与多馈公式 (77) &
\textbf{已实现}\quad \srcpath{compute\_mu} 对 $I_k''/I_r\le2$ 返回 1，并在 20/50/100/250 ms 曲线间插值；\srcpath{run\_short\_circuit\_detailed\_impl} 使用端电压跌落权重\\
稳态系数 $\lambda_{\max}/\lambda_{\min}$ &
\textbf{已实现}\quad \srcpath{Generator::sc\_lambda\_max/sc\_lambda\_min} 显式承载制造商/标准曲线值；缺值的单馈近端工况 fail closed；端子静态励磁取 $\lambda_{\min}$，多馈取 $I_{bmo}$\\
热等效 $m/n$ 标准曲线 &
\textbf{已实现}\quad \srcpath{calculate\_thermal\_factors\_sc}：附录 A 闭式、图 19 插值及全部极限处理\\
\bottomrule
\end{longtable}
\end{center}

\subsection{参考文献}
{\renewcommand{\section}[2]{}%
\begin{thebibliography}{9}\small
\bibitem{iec:iec60909-0} IEC 60909-0, \emph{Short-circuit currents in three-phase
a.c. systems -- Part 0: Calculation of currents}, 2016.
\bibitem{iec:iec60909-4} IEC TR 60909-4, \emph{Short-circuit currents in three-phase
a.c. systems -- Part 4: Examples for the calculation of short-circuit currents}, 2021.
\bibitem{iec:kundur} P.~Kundur, \emph{Power System Stability and Control},
McGraw-Hill, 1994.
\bibitem{iec:anderson} P.~M. Anderson, \emph{Analysis of Faulted Power Systems},
IEEE Press, 1995.
\end{thebibliography}}
